\documentclass{article}
\usepackage{amsmath}
\usepackage{graphicx}
\usepackage{booktabs}
\title{Il modello K-Nearest Neighbors}
\author{Prof. Fedeli Massimo}
\date{}
\begin{document}
	
	\maketitle
	
	Il modello \textbf{K-Nearest Neighbors} (KNN) è uno dei metodi più semplici dell'apprendimento automatico.
	
	L'idea di base è la seguente: per fare una previsione su un nuovo dato, il modello osserva i dati già noti che gli somigliano di più.
	
	
	\section*{Rappresentazione dei dati}
	
	Ogni oggetto (ad esempio uno studente) viene descritto tramite alcune caratteristiche dette \textit{feature}.
	
	Esempio:
	
	\begin{itemize}
		\item Ore di studio
		\item Numero di assenze
	\end{itemize}
	
	E' possibile rappresentare ogni studente come un vettore:
	
	\[
	\mathbf{x} = (x_1, x_2)
	\]
	
	dove
	
	\[
	x_1 = \text{ore di studio}, \quad x_2 = \text{assenze}
	\]
	
	Ogni studente diventa quindi un \textbf{punto nel piano cartesiano}.
	
	In generale si ha un insieme di dati noti (dataset):
	
	\[
	\mathcal{D} = \{(\mathbf{x}_1,y_1), (\mathbf{x}_2,y_2), \dots, (\mathbf{x}_n,y_n)\}
	\]
	
	dove
	
	\begin{itemize}
		\item $\mathbf{x}_i$ = caratteristiche (features)
		\item $y_i$ = risultato osservato (voto, esito, ecc.)  - target
	\end{itemize}
	
	\textbf{Esempio numerico.} Supponiamo di avere i seguenti dati su 5 studenti:
	
	\begin{center}
		\begin{tabular}{cccl}
			\toprule
			Studente & Ore di studio ($x_1$) & Assenze ($x_2$) & Esito ($y$) \\
			\midrule
			$A$ & 8 & 1 & Promosso \\
			$B$ & 2 & 5 & Bocciato  \\
			$C$ & 6 & 2 & Promosso  \\
			$D$ & 1 & 6 & Bocciato  \\
			$E$ & 7 & 0 & Promosso  \\
			\bottomrule
		\end{tabular}
	\end{center}
	
	Ipotizziamo che arrivi un nuovo studente:
	
	\[
	\mathbf{x}_{new} = (5,\; 3)
	\]
	
	di cui vogliamo prevedere il risultato.
	
	\section*{Misura di somiglianza}
	
	Per capire quali dati sono simili a $\mathbf{x}_{new}$ dobbiamo misurare la distanza.
	
	Applichiamo la formula della \textbf{distanza euclidea}:
	
	\[
	d(\mathbf{x}, \mathbf{y}) =
	\sqrt{
		\sum_{j=1}^{m}
		(x_j - y_j)^2
	}
	\]
	
	dove $m$ è il numero di caratteristiche.
	
	\textbf{Esempio numerico.} Calcoliamo la distanza tra il nuovo studente $\mathbf{x}_{new}=(5,3)$ e ciascuno dei cinque studenti noti:
	
	\[
	d(\mathbf{x}_{new}, A) = \sqrt{(5-8)^2 + (3-1)^2} = \sqrt{9+4} = \sqrt{13} \approx 3{,}61
	\]
	\[
	d(\mathbf{x}_{new}, B) = \sqrt{(5-2)^2 + (3-5)^2} = \sqrt{9+4} = \sqrt{13} \approx 3{,}61
	\]
	\[
	d(\mathbf{x}_{new}, C) = \sqrt{(5-6)^2 + (3-2)^2} = \sqrt{1+1} = \sqrt{2}  \approx 1{,}41
	\]
	\[
	d(\mathbf{x}_{new}, D) = \sqrt{(5-1)^2 + (3-6)^2} = \sqrt{16+9} = \sqrt{25} = 5{,}00
	\]
	\[
	d(\mathbf{x}_{new}, E) = \sqrt{(5-7)^2 + (3-0)^2} = \sqrt{4+9}  = \sqrt{13} \approx 3{,}61
	\]
	
	\section*{I K vicini più prossimi}
	
	Si selezionano i $K$ punti con distanza minima da $\mathbf{x}_{new}$.
	
	Questo insieme si indica con:
	
	\[
	N_K(\mathbf{x}_{new})
	\]
	
	cioè l'insieme dei $K$ vicini più prossimi.
	
	\textbf{Esempio numerico.} Ordiniamo gli studenti per distanza crescente:
	
	\begin{center}
		\begin{tabular}{ccc}
			\toprule
			Studente & Distanza & Esito \\
			\midrule
			$C$ & $1{,}41$ & Promosso \\
			$A$ & $3{,}61$ & Promosso \\
			$B$ & $3{,}61$ & Bocciato  \\
			$E$ & $3{,}61$ & Promosso  \\
			$D$ & $5{,}00$ & Bocciato  \\
			\bottomrule
		\end{tabular}
	\end{center}
	
	Con $K=3$ i vicini più prossimi sono $C$, $A$ e $B$, quindi:
	\[
	N_3(\mathbf{x}_{new}) = \{C,\, A,\, B\}
	\]
	
	\section*{Caso di classificazione}
	
	Se il problema è classificare (es: promosso / bocciato), il modello assegna la classe più frequente tra i vicini:
	
	\[
	\hat{y} =
	\text{mode}
	\{ y_i \mid \mathbf{x}_i \in N_K(\mathbf{x}_{new}) \}
	\]
	
	cioè la moda delle etichette dei vicini.
	
	\textbf{Definizione}. La moda è il valore  che si presenta con la massima frequenza all'interno di un insieme di dati
	
	\textbf{Esempio numerico.} Tra i 3 vicini $\{C, A, B\}$ le etichette sono:
	\[
	\{\text{Promosso},\; \text{Promosso},\; \text{Bocciato}\}
	\]
	La classe più frequente è \textit{Promosso} (2 voti su 3), quindi:
	\[
	\hat{y} = \text{Promosso}
	\]
	
	\section*{Caso di regressione}
	
	Se il problema è numerico (es: voto), il modello fa la media:
	
	\[
	\hat{y} =
	\frac{1}{K}
	\sum_{\mathbf{x}_i \in N_K(\mathbf{x}_{new})}
	y_i
	\]
	
	Quindi la previsione è la media dei valori dei vicini.
	
	\textbf{Esempio numerico.} Supponiamo ora che $y$ rappresenti il voto in decimi e che i dati siano:
	
	\begin{center}
		\begin{tabular}{ccc}
			\toprule
			Studente & $(x_1, x_2)$ & Voto ($y$) \\
			\midrule
			$C$ & $(6,\,2)$ & $8$ \\
			$A$ & $(8,\,1)$ & $9$ \\
			$B$ & $(2,\,5)$ & $5$ \\
			\bottomrule
		\end{tabular}
	\end{center}
	
	Con $K=3$ la previsione per il nuovo studente è:
	\[
	\hat{y} = \frac{1}{3}(8 + 9 + 5) = \frac{22}{3} \approx 7{,}33
	\]
	
	\section*{Il ruolo del parametro K}
	
	Il valore di $K$ è fondamentale.
	
	\begin{itemize}
		\item $K$ piccolo $\rightarrow$ decisioni molto locali
		\item $K$ grande $\rightarrow$ decisioni più stabili ma più ``medie''
	\end{itemize}
	
	\textbf{Esempio numerico.} Riprendendo il caso di classificazione con il nuovo studente $(5,3)$:
	
	\begin{itemize}
		\item Con $K=1$: unico vicino è $C$ (Promosso) $\;\Rightarrow\; \hat{y} = \text{Promosso}$
		\item Con $K=3$: vicini $C$, $A$, $B$ (2 Promosso, 1 Bocciato) $\;\Rightarrow\; \hat{y} = \text{Promosso}$
		\item Con $K=5$: tutti gli studenti (3 Promosso, 2 Bocciato) $\;\Rightarrow\; \hat{y} = \text{Promosso}$
	\end{itemize}
	
	In questo caso il risultato non cambia, ma con dataset più grandi e confini meno netti la scelta di $K$ può essere determinante.
	
	\section*{Riflessione}
	
	A differenza della regressione lineare, KNN non costruisce una funzione del tipo:
	
	\[
	y = f(x)
	\]
	
	ma utilizza direttamente i dati memorizzati.
	
	La previsione nasce dal confronto con esempi simili.
	
	\begin{center}
		Il modello non ricava una formula. Apprende per analogia.
	\end{center}
	
	\section*{Punti di forza}
	
	\begin{itemize}
		\item Facile da capire
		\item Non richiede ipotesi di linearità
		\item Si adatta a situazioni complesse
	\end{itemize}
	
	\section*{Limiti}
	
	\begin{itemize}
		\item Può essere lento con molti dati
		\item Sensibile alla scala delle variabili
		\item Funziona meno bene con molte dimensioni
	\end{itemize}
	
\end{document}